% ---------------------------------------------------------------------------
% Chapter 11: The Customs Form: Digitisation
% SKELETON (2026-11-15), laid out after Constitution.md rules B-E. Every
% \todo marks text, numbers or figures still to come.
% The scripts, records, logs and films of this stop will live in example/DIGI/,
% the technical companion is chapter/chapter_11_DIGI.md, and the CMSSW side
% (the DIGI step of the mixed event) is chapter 11 of JitJetSW, linked from
% example/DIGI/CMSSW/README.md.
% What this stop is: the SimHits of the family and the crowd become what
% the electronics write: charge, pulses sampled every 25 ns, pedestals and
% noise, ADC counts, zero suppression. Energy is added (noise, out-of-time
% pulses) and removed (thresholds). The checkpoints read the result.
% ---------------------------------------------------------------------------
%% chapter-local shorthands
\providecommand{\placeholderfig}[1]{\fbox{\parbox[c][4cm][c]{0.9\textwidth}{\centering\todo{#1}}}}
\providecommand{\bx}{\ensuremath{\mathrm{bx}}\xspace}
\providecommand{\ns}{\ensuremath{\,\mathrm{ns}}\xspace}
\providecommand{\Npe}{\ensuremath{N_{\mathrm{pe}}}\xspace}
\chapter{The Customs Form: Digitisation}
\chaptermark{The Customs Form}
\label{ch:customs_form}

\begin{journeybox}
The weighing is over, and now it has to be written down. The customs form
has boxes of fixed size, its own units and its own rounding; a small item
is not worth a line and is left off, the clerk's hand trembles a little on
every entry, and the ink of the previous traveller has not dried on the
page. What travels on is not the luggage but the form. \todo{Polish the
metaphor paragraph (B1).}
\end{journeybox}

\physicsline{from deposited energy to charge: light yield, photostatistics,
gains and saturation; the pulse and the $25\ns$ clock, and the out-of-time
crowd inside a pulse; pedestals, electronic noise and quantisation into ADC
counts; zero suppression and selective readout; the tracker's charge
sharing and thresholds. The one number to take away: \todo{the energy of
our jet in digitised cells above threshold, against the SimHits of
\cref{ch:crowd_at_the_border}}.}

\section{What we will see at this stop}
\label{sec:customs_form:what_we_compute}

% (C1-1) motivation, connected to the previous chapter
\todo{Motivation (C1-1). \Cref{ch:crowd_at_the_border} ended with energy in
the active volumes, the family's and the crowd's added together. No
detector reads energy: it reads light or charge, shaped into a pulse,
sampled by the clock and written in counts. This stop is the conversion,
and the first place where the event is what the real detector writes.
Give the four numbers to take away.}

\paragraph{The film.}
\label{sec:customs_form:animation}
% (C1-2) one paragraph, at most one figure of key frames
\todo{One paragraph on the eleventh film of the playlist of \jitjetworld.}

\begin{figure}[htbp]\centering
\placeholderfig{Key frames of the film, from
\file{example/DIGI/Standalone/animate_standalone_DIGI_stills.pdf}.}
\caption[Key frames of the animated digitisation stage]{Key frames of the
animated digitisation stage. \todo{Caption.}}
\label{fig:customs_form:stills}
\end{figure}

\paragraph{The picture to carry: one cell, from energy to counts.}
% (C1-3) a sketch or diagram
\todo{The picture to carry: one ECAL crystal and one HCAL tower of our jet,
the energy turned into photoelectrons, the pulse with the out-of-time
pulses under it, the ten samples, the pedestal and the counts written.
TikZ in \file{chapter/tikz/digi_*.tex}.}

\begin{figure}[htbp]\centering
\placeholderfig{TikZ: one cell from energy to ADC counts
(\file{chapter/tikz/digi_pulse.tex}).}
\caption[The customs form]{The customs form. \todo{Caption.}}
\label{fig:customs_form:pulse_sketch}
\end{figure}

\paragraph{The number, three ways.}
\label{sec:customs_form:validation}
% (C1-4) the agreement table, one or two paragraphs
\todo{One or two paragraphs on the agreement table: the standalone script
(the digitisers ported to NumPy), the official package (the digitiser
models with the 2024 conditions), and \CMSSW (the DIGI step of JitJetSW).
How each column is produced is in the companion,}
\begin{center}
  \companion{chapter_11_DIGI}
\end{center}

\begin{table}[htbp]\centering\small
\caption[The customs form obtained three ways]{The customs form obtained
three ways. \todo{Caption: what each column is.}}
\label{tab:customs_form:numbers}
\begin{tabular}{lrrr}\toprule
 & standalone & official & \CMSSW 2024 \\ \midrule
photoelectrons per GeV: ECAL / HCAL & \todo{} & \todo{} & \todo{} \\
noise per channel [MeV] & \todo{} & \todo{} & \todo{} \\
out-of-time energy in our jet's pulses [GeV] & \todo{} & \todo{} & \todo{} \\
channels read out after zero suppression & \todo{} & \todo{} & \todo{} \\
energy lost below threshold in our jet [GeV] & \todo{} & \todo{} & \todo{} \\
tracker clusters in our jet & \todo{} & \todo{} & \todo{} \\
\bottomrule
\end{tabular}
\end{table}

\section{Let us dive into the details: the theory}
\label{sec:customs_form:theory}
% Every subsection: (C2-1) motivation, (C2-2) TikZ sketch, (C2-3) equation,
% (C2-4) numbers for our event from the -v log, inline, (C2-5) one plot.

\subsection{From energy to charge: light, photoelectrons and gains}
\label{sec:customs_form:charge}
\todo{(C2-1) The crystal makes light, the scintillator makes light, the
silicon makes charge; each conversion counts quanta and fluctuates.
(C2-2) TikZ: crystal, photodetector and preamplifier.
(C2-3) $\Npe=E\cdot Y$ and the photostatistics term $\sigma/E=1/\sqrt{\Npe}$;
the saturation of a SiPM.
(C2-4) Numbers from the \code{-v} log.
(C2-5) Plot: the photostatistics term against energy.}

\subsection{The pulse and the clock}
\label{sec:customs_form:pulse}
\todo{(C2-1) A hit is a pulse that lasts longer than $25\ns$, and the pulses
of the crossings before ours sit under it.
(C2-2) TikZ: the pulse and its samples, with the out-of-time pulses.
(C2-3) The pulse shape $f(t)$ and the sampled amplitude
$S_i=\sum_{\bx}A_{\bx}\,f(t_i-25\ns\cdot\bx)$.
(C2-4) Numbers from the \code{-v} log.
(C2-5) Plot: the ten samples of the hottest crystal of our jet.}

\subsection{Pedestals, noise and ADC counts}
\label{sec:customs_form:noise}
\todo{(C2-1) The form has boxes of fixed size and the clerk's hand trembles.
(C2-2) TikZ: the pedestal, the noise band and the ADC steps.
(C2-3) The quantisation noise $q/\sqrt{12}$, the electronic noise and the
gain switching.
(C2-4) Numbers from the \code{-v} log.
(C2-5) Plot: the noise of our jet's channels.}

\subsection{What is not written: zero suppression and selective readout}
\label{sec:customs_form:zero_suppression}
\todo{(C2-1) Most channels hold noise, and writing them all is impossible.
(C2-2) TikZ: the selective readout regions around our jet.
(C2-3) The threshold $n\sigma_{\mathrm{noise}}$ and the energy lost below it.
(C2-4) Numbers from the \code{-v} log.
(C2-5) Plot: the energy lost against the threshold.}

\subsection{The tracker: charge sharing and clusters}
\label{sec:customs_form:tracker}
\todo{(C2-1) A charged member crossing silicon leaves charge on several
pixels or strips.
(C2-2) TikZ: drift, Lorentz angle and charge sharing.
(C2-3) The Lorentz drift $\tan\theta_L=\mu_H B$ and the cluster position.
(C2-4) Numbers from the \code{-v} log.
(C2-5) Plot: the cluster sizes of our jet's tracks.}

\paragraph{The recipe.}
\todo{The recipe, step by step, that the standalone script runs.}

\section{The energy ledger: our jet at this stage}
\label{sec:customs_form:ledger}
% (C3-1) the ledger table, "The same event in CMSSW", one figure page
\todo{What the record \file{example/DIGI/event_DIGI.*} holds, and what the
ledger says.}

\begin{ledgerboth}{The Customs Form}{tab:customs_form:ledger}{%
\todo{What the two blocks are.}}
\ledgerblock{Standalone (seed 1)}
\ledgerrow{the photograph, \code{ak4GenJets} jet [0] (chapter 7)}{99.83}{14.83}{31}{--}
\ledgerrow{with the crowd's SimHits (chapter 10)}{\todo{}}{--}{\todo{}}{\todo{}}
\ledgerrow{digitised, above threshold}{\todo{}}{--}{\todo{}}{\todo{}}
\ledgerblock{\CMSSW (run 1, event 1)}
\ledgerrow{the photograph, \code{ak4GenJets} jet [0] (chapter 7)}{117.79}{16.72}{24}{--}
\ledgerrow{digitised, above threshold}{\todo{}}{--}{\todo{}}{\todo{}}
\end{ledgerboth}

\paragraph{The same event in \CMSSW.}
\label{sec:customs_form:cmssw_event}
\todo{The DIGI step of JitJetSW: the digis of our jet, read from the log of
the analyser.}

\begin{figure}[p]\centering
\placeholderfig{Standalone display: \file{DIGI/journey_DIGI.pdf}.}\\[2mm]
\placeholderfig{\CMSSW display: \file{DIGI/CMSSW/digiAnalyzer_11_DIGI.pdf}.}
% when the files exist:
% \ledgerdisplays
%   {DIGI/journey_DIGI.pdf}{Standalone: ...}{fig:customs_form:journey}
%   {DIGI/CMSSW/digiAnalyzer_11_DIGI.pdf}{\CMSSW: ...}{fig:customs_form:cmssw}
\caption[The digitised event, standalone and \CMSSW]{The digitised event,
standalone (top) and \CMSSW (bottom). \todo{Caption.}}
\label{fig:customs_form:displays}
\end{figure}

\subsection{What this stop is, and what it is not}
\label{sec:customs_form:caveats}
\begin{itemize}
  \item \todo{Digis are counts, not energies: turning them back into
    energy is the local reconstruction of \cref{ch:footprints}.}
  \item \todo{The conditions are the 2024 ones, one set for the whole year.}
  \item \todo{One event is one sample.}
\end{itemize}

\section{Further reading}
\label{sec:customs_form:further_reading}
The signal amplitudes of the ECAL and their reconstruction are in
\cite{CMS:2020uim}, its calibration and resolution in \cite{CMS:2013lxn};
the Run-3 electronics of the calorimeters, including the SiPM readout of the
HCAL, are in \cite{CMS:2024aqx}, and calorimetry in general in
\cite{Wigmans:2000vf}. \todo{Literature survey: the HCAL and tracker
digitisers, selective readout.}

\begin{pitfalls}
\begin{itemize}[nosep]
  \item \todo{Pitfall 1.}
  \item \todo{Pitfall 2.}
\end{itemize}
\end{pitfalls}

\begin{exercises}
\item \todo{Exercise 1.}
\item \todo{Exercise 2.}
\end{exercises}

\subsection*{Open problems in the context of Phase-2 LHC}
\label{sec:customs_form:open_problems}
\todo{Literature survey: the Phase-2 front ends, the high-granularity
calorimeter~\cite{CMS:2017jpq}, the Phase-2 tracker~\cite{CMS:2017lum}
and timing~\cite{Butler:2019rpu}.}
\begin{enumerate}
  \item \todo{Open problem 1.}
  \item \todo{Open problem 2.}
\end{enumerate}
The next stop is the first checkpoint, where the form is read in a hurry.
